\section{An integral inverse gate and the residual cost of cancellation}
\label{il-section}

The rational models of Section~\ref{rd-section} require an integral
inverse condition. We give both a sufficient local domain and a sharp
cost of cancellation for any residual at the prescribed exponent.
Neither supplies a point-height upper bound.

Let $\mathcal O=\mathbb Z[\zeta]$, $\zeta^2-\zeta+1=0$ and
$\gamma=1+\zeta$. For an integer $n\ge2$, suppose
\begin{equation}\label{il-input}
 \tau=u\gamma^e v,\quad e\in\{0,1\},\quad V=\operatorname N v,
 \qquad w\text{ primitive unramified},\quad R=\operatorname N w>1,
\end{equation}
where $u$ is a unit and $v$ is primitive oriented unramified.
Then $\tau$ is primitive, $\operatorname N\tau=3^eV$ and
$3\nmid VR$. Put
\begin{equation}\label{il-reduced}
 z=\tau w^n=A+B\zeta,\quad C=\gcd(|A|,|B|),\quad
 M_{\rm red}=\operatorname N(z/C)=\frac{3^eVR^n}{C^2}.
\end{equation}
Here $C>0$. The reduced pair is primitive, without a positivity
assumption in the following single-profile statements.

\begin{lemma}[The exact opposite-orientation content]\label{il-content}
The integer $C$ is prime to three and divides $V$. Its prime divisors
are exactly the split rational primes where $\tau$ and $w$ have
opposite orientations. At such a prime $p$, put
$a_p=v_p(V)>0$ and $f_p=v_p(R)>0$. Then
\begin{equation}\label{il-opposite-depths}
 v_p(C)=\min(a_p,nf_p),\qquad
 v_p(M_{\rm red})=|a_p-nf_p|.
\end{equation}
A shared prime with the same orientation contributes no content.
\end{lemma}
\begin{proof}
Primitivity in the Eisenstein UFD means that at a split rational
prime each input contains at most one orientation. No inert prime
occurs in either primitive norm. At the ramified prime $\gamma$,
the root is a unit and the multiplier has depth $e\le1$.
Rational divisibility by three would require depth at least two,
because $\gamma^2=3\zeta$. Thus $3\nmid C$.

At a split $p$, divisibility by the rational prime is divisibility
by both conjugate Eisenstein primes. Equal orientations retain only
one factor. Opposite orientations give depths $a_p$ and $nf_p$,
so the rational content has their minimum as depth. After division,
the norm exponent is their absolute difference. This proves
\eqref{il-opposite-depths} and the support assertion. Finally
$\min(a_p,nf_p)\le a_p$, so $C\mid V$.
\end{proof}

The more general content bound of Lemma~\ref{pc-content} follows
from an adjugate and Bezout argument. The exact minimum above uses
primitive oriented factors; primitivity of $\tau$ is material to
the next result.

\begin{theorem}[A bill for any output residual]\label{il-residual-bill}
Suppose the reduced norm has any representation
\begin{equation}\label{il-output-profile}
 M_{\rm red}=3^eV'(R')^n,\qquad V',R'\in\mathbb Z_{>0}.
\end{equation}
The root $R'$ may differ from $R$ and may be one. Then
\begin{equation}\label{il-radical-bill}
 \operatorname{rad}(C)^n\mid VV',
\end{equation}
where $\operatorname{rad}(1)=1$. In particular $C>1$ implies
\begin{equation}\label{il-log-threshold}
 \frac{\log V+\log V'}n\ge\log7.
\end{equation}
\end{theorem}
\begin{proof}
Every prime $p\mid C$ is split, so $p\ge7$. If $a_p\ge n$,
$V$ alone contains $p^n$. If $1\le a_p<n$, then
$nf_p\ge n>a_p$, and \eqref{il-opposite-depths} gives
$v_p(M_{\rm red})=nf_p-a_p$. The nonnegative valuation $v_p(V')$
in \eqref{il-output-profile} is congruent to $-a_p$ modulo $n$,
so it is at least $n-a_p$. Thus in both cases
\[
 v_p(V)+v_p(V')\ge n.
\]
Multiplication over the distinct primes dividing $C$ proves
\eqref{il-radical-bill}. Its logarithm gives
$(\log V+\log V')/n\ge\log\operatorname{rad}(C)$, which is
at least $\log7$ when $C>1$.
\end{proof}

A canonical $n$-free representation \eqref{il-output-profile}
always exists by taking the norm exponents modulo $n$ away from
three. The exponent at three remains exactly $e$ by
\eqref{il-reduced}. Thus the theorem bounds the smallest possible
output residual and continues to hold when the numerical root is
changed. Its threshold is sharp: the family of
Theorem~\ref{pc-inverse-boundary} has
\[
 C=7,\qquad V=7,\qquad V'=7^{n-1},
\]
so equality holds in \eqref{il-radical-bill} and
\eqref{il-log-threshold}. This is a single-map example, not a
rational point asserted to lie on the full double fiber product.

\begin{theorem}[An explicit integral inverse domain]\label{il-double-inverse}
Fix integers $h,g\ge2$ and $\tau_0=u_0\gamma^e v_0$,
$\tau_1=u_1v_1$ as in
Section~\ref{rd-section}, with primitive oriented unramified
$v_i$ and $V_i=\operatorname N v_i$. Let $(t_0,t_1)$ be a
rational point of the full fiber product, with
\begin{equation}\label{il-point-domain}
 x=L_{\tau_0,h}(t_0)>0\text{ finite},\qquad
 L_{\tau_1,g}(t_1)=\frac{x}{x^2+x+1}.
\end{equation}
Choose primitive integral homogeneous coordinates for $t_i$ and
write $w_i=r_i+s_i\zeta$. Assume these roots are unramified
nonunits. Put
\[
 R=\operatorname N w_0,\quad Q=\operatorname N w_1,\qquad
 C_0=\operatorname{cont}(\tau_0w_0^h),\quad
 C_1=\operatorname{cont}(\tau_1w_1^g).
\]
If $x=a/b$ in lowest positive terms, then
\begin{equation}\label{il-two-norms}
 a^2+ab+b^2=\frac{3^eV_0R^h}{C_0^2},\qquad
 F(a,b)=\frac{V_1Q^g}{C_1^2}.
\end{equation}
In particular $C_0=C_1=1$ gives actual oriented profiles with
the original residuals and exponents, up to unit signs. More
generally, numerical profiles with the same $R,Q,h,g$ and
positive integral residuals at most $V_0,V_1$ exist exactly when
\begin{equation}\label{il-square-content-condition}
 C_0^2\mid V_0,\qquad C_1^2\mid V_1.
\end{equation}
\end{theorem}
\begin{proof}
The first output has ratio $a/b$. Dividing its coordinates by
their content and changing the common sign gives exactly $(a,b)$.
Taking norms proves the first equation of \eqref{il-two-norms}.
Put $M=a^2+ab+b^2$. Since $\gcd(ab,M)=1$, the second output,
whose ratio is $ab/M>0$, reduces in the same way to $(ab,M)$.
Its norm is $F(a,b)$, proving the other equation. All actual
denominators are nonzero since $a,b,M>0$. A coordinate $s_i=0$
would make the primitive root a unit, already excluded here.

When both contents are one these are the actual oriented
factorizations, with at most sign changes to their units. For
general content, \eqref{il-two-norms} determines the residuals
with the same numerical root norms uniquely as $V_i/C_i^2$.
They are positive integers precisely under
\eqref{il-square-content-condition}, and then are at most $V_i$.
The numerical profiles have actual oriented liftings by the same
UFD argument as Lemma~\ref{dc-second-lift}. The original root
orientations need not be preserved after division. Nor does failure
of \eqref{il-square-content-condition} exclude a profile with a
different numerical root norm.
\end{proof}

\paragraph{The finite local condition in a low-residual domain.}
On either side, write $n$ for its exponent, $V$ for the initial
residual, and $V'$ for any desired output residual at that exponent.
If $(\log V+\log V')/n<\log7$, Theorem~\ref{il-residual-bill}
forces content one, even if the numerical root is changed.

Content one has a finite local description. For each $p\mid V$,
the primitive multiplier is divisible by one split prime factor.
Its product with $w^n$ has no content at $p$ exactly when $w$ is
not divisible by the opposite factor. In primitive coordinates
$[r:s]$, this excludes the single point $r/s=-\theta$ of
$\mathbb P^1(\mathbb F_p)$, where $\theta$ is the image of
$\zeta$ for the opposite factor. A norm-zero primitive residue
has $s\ne0$, so no point at infinity is omitted. No other prime
gives a content condition; unramifiedness of the root at three
is the separate hypothesis already stated.

These exclusions on both sides suffice to reconstruct the original
integral residuals $V_0,V_1$. They do not guarantee an arbitrarily
specified smaller output-residual budget. For such a smaller
target, the assertion is the necessary exclusion of nonzero content
under the joint input/output bound. The statement concerns the
particular rational point and its root coordinates, not uniqueness
of every profile of the resulting seed.

The bill \eqref{il-radical-bill} and the integral inverse domain
do not prove the existence of suitable points or bound their heights.
They identify the arithmetic domain on which that remaining question
must be addressed. These are ordinary proofs; no full Lean
formalization of the oriented or geometric assertions is claimed.
